\documentclass[10pt,a4paper]{amsart}
\usepackage{amsmath,amssymb,amsthm,graphicx}
\usepackage{verbatim}
\newenvironment{reference}{\comment}{\endcomment}
%\newenvironment{reference}{}{}
\newenvironment{slogan}{\comment}{\endcomment}
\newenvironment{history}{\comment}{\endcomment}

% For commutative diagrams we use Xy-pic
\usepackage[all]{xy}

% We use 2cell for 2-commutative diagrams.
\xyoption{2cell}
% Theorem environments.
%
\theoremstyle{plain}
\newtheorem{theorem}[subsection]{Theorem}
\newtheorem{proposition}[subsection]{Proposition}
\newtheorem{lemma}[subsection]{Lemma}

\theoremstyle{definition}
\newtheorem{definition}[subsection]{Definition}
\newtheorem{example}[subsection]{Example}
\newtheorem{exercise}[subsection]{Exercise}
\newtheorem{situation}[subsection]{Situation}

\theoremstyle{remark}
\newtheorem{remark}[subsection]{Remark}
\newtheorem{remarks}[subsection]{Remarks}

\numberwithin{equation}{subsection}

% Macros
%
\def\lim{\mathop{\mathrm{lim}}\nolimits}
\def\colim{\mathop{\mathrm{colim}}\nolimits}
\def\Spec{\mathop{\mathrm{Spec}}}
\def\Hom{\mathop{\mathrm{Hom}}\nolimits}
\def\Ext{\mathop{\mathrm{Ext}}\nolimits}
\def\SheafHom{\mathop{\mathcal{H}\!\mathit{om}}\nolimits}
\def\SheafExt{\mathop{\mathcal{E}\!\mathit{xt}}\nolimits}
\def\Sch{\mathit{Sch}}
\def\Mor{\mathop{\mathrm{Mor}}\nolimits}
\def\Ob{\mathop{\mathrm{Ob}}\nolimits}
\def\Sh{\mathop{\mathit{Sh}}\nolimits}
\def\NL{\mathop{N\!L}\nolimits}
\def\CH{\mathop{\mathrm{CH}}\nolimits}
\def\proetale{{pro\text{-}\acute{e}tale}}
\def\etale{{\acute{e}tale}}
\def\QCoh{\mathit{QCoh}}
\def\Ker{\mathop{\mathrm{Ker}}}
\def\Im{\mathop{\mathrm{Im}}}
\def\Coker{\mathop{\mathrm{Coker}}}
\def\Coim{\mathop{\mathrm{Coim}}}

% Boxtimes
%
\DeclareMathSymbol{\boxtimes}{\mathbin}{AMSa}{"02}

%
% Macros for moduli stacks/spaces
%
\def\QCohstack{\mathcal{QC}\!\mathit{oh}}
\def\Cohstack{\mathcal{C}\!\mathit{oh}}
\def\Spacesstack{\mathcal{S}\!\mathit{paces}}
\def\Quotfunctor{\mathrm{Quot}}
\def\Hilbfunctor{\mathrm{Hilb}}
\def\Curvesstack{\mathcal{C}\!\mathit{urves}}
\def\Polarizedstack{\mathcal{P}\!\mathit{olarized}}
\def\Complexesstack{\mathcal{C}\!\mathit{omplexes}}
% \Pic is the operator that assigns to X its picard group, usage \Pic(X)
% \Picardstack_{X/B} denotes the Picard stack of X over B
% \Picardfunctor_{X/B} denotes the Picard functor of X over B
\def\Pic{\mathop{\mathrm{Pic}}\nolimits}
\def\Picardstack{\mathcal{P}\!\mathit{ic}}
\def\Picardfunctor{\mathrm{Pic}}
\def\Deformationcategory{\mathcal{D}\!\mathit{ef}}
\usepackage{cleveref}
\crefname{originalreference}{original reference}{original references}
\Crefname{originalreference}{Original reference}{Original references}
\makeatletter
\newlabel{lemma-take-pth-root}{{09HF}{source}{}{sourcepointer-1}{}}
\newlabel{lemma-take-pth-root@cref}{{[originalreference][0][]09HF}{[1][0][]source}{}{}{}}
\makeatother
\title{p-th roots in separable extensions: an element whose minimal polynomial has p-th-power coefficients is a p-th power}
\author{The Stacks Project Authors}
\date{stacks-project:fields.tex:local-snapshot-20261012; GNU FDL 1.2 or later}
\begin{document}
\maketitle
\noindent\textit{Editorial note.} All delivered mathematical text is the author's own, copied from the pinned Stacks Project chapter file fields.tex; no statement or proof step was rewritten, reordered or removed.\par\medskip
\noindent\textit{Editorial note.} The counted proof is the author's complete argument: The statement's hypothesis that L/K is separable enters only through the distinct-roots clause of (2b); the lemma feeds the p-basis machinery later in the chapter.\par\medskip
\noindent\textit{Editorial note.} The label lemma-pth-root names the counted statement, and the proof environment immediately following it in the source is delivered verbatim as its proof.\par\medskip
\noindent\textit{Editorial note.} The counted statement is carried by the author's own lemma environment in the chapter source.\par\medskip
\noindent\textit{Allowed standard prerequisites (editorial).}\begin{itemize}
\item Same chapter: lemma-take-pth-root (if alpha is not a p-th power in L then T\textasciicircum{}p - alpha is irreducible over L); separability of L/K gives the minimal polynomial of alpha pairwise distinct roots
\end{itemize}\medskip
\ifdefined\hypertarget\hypertarget{sourcepointer-1}{}\fi
\noindent\textit{Source reference.} Original reference 09HF: sources/stacks-fields/pinned-original-source.tex, source lines 1834--1834. Prerequisite pointer invoked at 3985; it names a standard fact of the same chapter that the counted argument uses, and every step of the counted proof built on it is retained verbatim in the delivered body.\par\medskip
\begin{lemma}
\label{lemma-pth-root}
Let $K$ be a field of characteristic $p > 0$. Let $L/K$ be a separable
algebraic extension. Let $\alpha \in L$.
\begin{enumerate}
\item If the coefficients of the minimal polynomial of $\alpha$
over $K$ are $p$th powers in $K$ then $\alpha$ is a $p$th
power in $L$.
\item More generally, if $P \in K[T]$ is a polynomial such that (a) $\alpha$
is a root of $P$, (b) $P$ has pairwise distinct roots in an algebraic closure,
and (c) all coefficients of $P$ are $p$th powers, then $\alpha$ is a
$p$th power in $L$.
\end{enumerate}
\end{lemma}

\begin{proof}
It follows from the definitions that (2) implies (1). Assume $P$ is as in (2).
Write $P(T) = \sum\nolimits_{i = 0}^d a_i T^{d - i}$ and $a_i = b_i^p$.
The polynomial $Q(T) = \sum\nolimits_{i = 0}^d b_i T^{d - i}$ has distinct
roots in an algebraic closure as well, because the roots of $Q$
are the $p$th roots of the roots of $P$. If $\alpha$ is not a $p$th power,
then $T^p - \alpha$ is an irreducible polynomial over $L$
(Lemma \ref{lemma-take-pth-root}).
Moreover $Q$ and $T^p - \alpha$ have a root in common in
an algebraic closure $\overline{L}$.
Thus $Q$ and $T^p - \alpha$ are not relatively prime, which
implies $T^p - \alpha | Q$ in $L[T]$. This contradicts the fact that
the roots of $Q$ are pairwise distinct.
\end{proof}

\end{document}
